% !TEX root = ./main.tex
\chapter{Studio empirico}
\label{chap:studio-empirico}

Dopo la valutazione tecnica del prototipo, restava da capire come SOULFRAME reggesse quando passava davvero in mano alle persone. In questa parte della tesi descrivo il test empirico svolto sul prototipo, con un confronto tra una condizione neutra e una condizione personale. L'obiettivo non era dimostrare qualcosa in astratto, ma capire se la personalizzazione dell'avatar cambiasse davvero il modo in cui l'agente veniva percepito e vissuto durante la sessione.

Lo studio è stato impostato come \textit{between-subjects}\footnote{Ogni partecipante ha provato una sola condizione. Così il test restava più semplice da gestire e la seconda prova non veniva influenzata dalla prima.}: ogni partecipante ha svolto pre-test, interazione e post-test in una sola sessione, senza passare anche per l'altra condizione. La prova serviva a capire se un avatar più vicino alla persona, per voce e memoria iniziale, risultasse davvero più credibile, oppure se finisse soprattutto per rendere più visibili i limiti del sistema.

\section{Metodo}
\label{sec:empirical-method}

\paragraph{Partecipanti.}
\label{subsec:empirical-participants}

Il campione complessivo è formato da 30 partecipanti, reclutati in modo informale e coinvolti in sessioni individuali svolte in presenza. Nel capitolo mantengo la lettura del test come confronto tra due gruppi da 15 partecipanti ciascuno, perché questa era l'impostazione con cui la prova era stata pensata e presentata. Tutti i dati sono stati raccolti in forma anonima e ogni partecipante ha dato il consenso alla prova prima dell'inizio della sessione.

Il gruppo risulta abbastanza eterogeneo per età e titolo di studio. L'età media complessiva è pari a 33.2 anni, con mediana 25.5 e IQR\footnote{\textit{Interquartile Range}: l'intervallo compreso tra primo e terzo quartile, utile qui per dare un'idea della dispersione senza farsi guidare troppo dai valori estremi.} 24--40. Dal punto di vista del genere, i partecipanti sono 17 maschi e 13 femmine. Per il titolo di studio prevalgono scuola superiore (16 casi), laurea triennale (9) e scuola media (5). La Tabella~\ref{tab:empirical-sample} riassume questi dati.

\begin{table}[H]
\centering
\small
\begin{tabularx}{\textwidth}{lX}
\hline
\textbf{Indicatore} & \textbf{Valore osservato} \\
\hline
Partecipanti complessivi & 30 \\
Disegno sperimentale & due condizioni tra gruppi (neutra / personale) \\
Età & media 33.2 anni; mediana 25.5; IQR 24--40 \\
Genere & 17 maschi, 13 femmine \\
Titolo di studio & 16 scuola superiore, 9 laurea triennale, 5 scuola media \\
Modalità di reclutamento & convenienza / contatto diretto \\
Setting & sessioni individuali su laptop Windows \\
\hline
\end{tabularx}
\caption{Descrizione sintetica del campione utilizzato nello studio empirico.}
\label{tab:empirical-sample}
\end{table}

\paragraph{Materiali dello studio.}
\label{subsec:empirical-materials}

Per il test è stata usata la build Windows del prototipo, eseguita sullo stesso laptop già descritto nel capitolo tecnico. Ho preferito questa configurazione perché, nelle prove precedenti, risultava più regolare del WebGL soprattutto sul piano del playback audio e della continuità generale della sessione. Tutti i partecipanti hanno interagito solo tramite tastiera e non hanno avuto accesso alle funzioni di modifica voce e modifica memoria durante la prova.

Nel prototipo era presente una \textit{test mode} dedicata, richiamabile dal menu principale con la sequenza \texttt{T-E-S-T}. Questa modalità serviva a standardizzare il più possibile il test: separava gli avatar generici da quelli personali, disabilitava i pulsanti di modifica in \texttt{MainMode}, forzava il setup memoria in cinque passaggi per la condizione personale e faceva ripartire ogni prova da una nuova sessione conversazionale.

La frase usata nel setup voce era la seguente:

\begin{quote}
``Mi chiamo \_\_\_\_. Oggi è il giorno \{0\} del mese di \{1\} dell'anno \{2\}. Il cielo oggi è \_\_\_\_\_. Sto facendo una prova di voce calma e chiara, quindi leggo lentamente questa frase guidata con attenzione, regolarmente. Ultime parole semplici per completare il test.''
\end{quote}

Nel passaggio reale i tre placeholder numerici venivano sostituiti con giorno, mese e anno correnti, mentre gli spazi lasciati vuoti venivano completati dal partecipante durante la registrazione. Nella condizione personale, inoltre, il setup memoria non avveniva con un testo libero unico, ma con cinque campi guidati:

\begin{table}[H]
\centering
\small
\begin{tabularx}{0.92\textwidth}{>{\raggedright\arraybackslash}p{0.30\textwidth}X>{\raggedleft\arraybackslash}p{0.12\textwidth}}
\hline
\textbf{Campo} & \textbf{Indicazione data al partecipante} & \textbf{Max caratteri} \\
\hline
Percorso di studi & corso di laurea, università, anno o anno di conclusione & 200 \\
Lavoro / occupazione & lavoro attuale o principale occupazione & 150 \\
Hobbies e interessi & due o tre attività del tempo libero praticate regolarmente & 200 \\
Città / contesto di vita & città di residenza e contesto generale & 150 \\
Tratti personali & due o tre aggettivi con cui descriversi & 100 \\
\hline
\end{tabularx}
\caption{Struttura del setup memoria guidato nella condizione personale.}
\label{tab:empirical-memory-setup}
\end{table}

Le due condizioni sono state organizzate così:

\begin{table}[H]
\centering
\small
\begin{tabularx}{\textwidth}{>{\raggedright\arraybackslash}p{0.18\textwidth}X X}
\hline
\textbf{Condizione} & \textbf{Configurazione dell'avatar} & \textbf{Materiale dato al sistema} \\
\hline
Neutra & avatar generico standardizzato, voce standardizzata, memoria di base non personalizzata; nuova sessione conversazionale per ogni prova & nessuna biografia personale del partecipante inserita nella memoria iniziale \\
Personale & avatar personale, frase di setup voce standardizzata, memoria guidata in cinque campi con limiti di caratteri; nuova sessione conversazionale per ogni prova & profilo personale del partecipante costruito tramite i cinque campi memoria \\
\hline
\end{tabularx}
\caption{Configurazione delle due condizioni sperimentali.}
\label{tab:empirical-conditions}
\end{table}

Nella condizione neutra tutti i partecipanti partivano dalla stessa configurazione di base dell'avatar. Allo stesso tempo, le prove non condividevano la stessa conversazione attiva, così ogni sessione veniva osservata come episodio separato.

Per la raccolta dati ho usato quattro strumenti standardizzati. La Rosenberg Self-Esteem Scale (RSES) \cite{Rosenberg1965} è stata somministrata prima della sessione. Il PANAS \cite{Watson1988PANAS} è stato proposto sia nel pre-test sia nel post-test, così da osservare la variazione dell'affettività positiva e negativa. Nel post-test ho poi usato il Godspeed Questionnaire Series \cite{Bartneck2009Godspeed} per descrivere come veniva percepito l'avatar, insieme a una versione adattata del Multidimensional Measure of Trust (MDMT) \cite{MalleUllman2021MDMT} per la fiducia nell'agente.

Nella fase di analisi ho trattato i questionari in modo piuttosto diretto. Per la RSES ho ricodificato gli item inversi e calcolato un punteggio totale. Per il PANAS ho separato la componente positiva da quella negativa e ho letto il cambiamento come differenza tra post-test e pre-test. Per la resa dell'avatar ho mantenuto le cinque dimensioni previste dal Godspeed --- antropomorfismo, animazione, gradevolezza, intelligenza percepita e sicurezza percepita --- anche se nel protocollo testuale compare per errore la dicitura ``quattro dimensioni''. Per l'MDMT, invece, ho preferito un indice medio complessivo su scala 0--7, così da avere una lettura più compatta della fiducia.

I punteggi dei questionari sono stati calcolati nel modo seguente. Per la RSES ho invertito gli item 2, 5, 6, 8 e 9 e poi sommato i dieci valori:
\[
x_i^*=
\begin{cases}
x_i & \text{se l'item non è invertito}\\
5-x_i & \text{se l'item è invertito}
\end{cases}
\qquad
RSES_{tot}=\sum_{i=1}^{10}x_i^*
\]

Per il PANAS ho separato i cinque item positivi dai cinque negativi, sia prima sia dopo la sessione:
\[
PA_{pre}=\sum_{i=1}^{5}x_i^{pre}, \qquad NA_{pre}=\sum_{i=6}^{10}x_i^{pre}
\]
\[
PA_{post}=\sum_{i=1}^{5}x_i^{post}, \qquad NA_{post}=\sum_{i=6}^{10}x_i^{post}
\]
\[
\Delta PA = PA_{post}-PA_{pre}, \qquad \Delta NA = NA_{post}-NA_{pre}
\]
In questa lettura, un valore \(\Delta PA > 0\) indica un aumento dell'affettività positiva, mentre \(\Delta NA < 0\) indica una riduzione dell'affettività negativa e quindi un andamento più favorevole.

Per il Godspeed ho lavorato sulla media di ciascuna dimensione:
\[
G_k=\frac{1}{m_k}\sum_{j=1}^{m_k}x_j
\]
e, nel caso della sicurezza percepita, ho riportato sullo stesso orientamento gli item inversi con:
\[
x_j^*=6-x_j
\]

Infine, per l'MDMT ho usato un indice medio complessivo:
\[
Trust=\frac{1}{20}\sum_{i=1}^{20}x_i
\]

\paragraph{Procedura.}
\label{subsec:empirical-procedure}

\vspace*{\fill}
La sessione seguiva un flusso abbastanza lineare. Dopo una breve introduzione, ogni partecipante compilava il pre-test con RSES e PANAS. A questo punto iniziava la parte operativa: nella condizione neutra l'avatar era già pronto, mentre nella condizione personale venivano completati il setup voce e il setup memoria guidato. Solo dopo questo passaggio iniziava la conversazione vera e propria con l'avatar. Al termine della prova venivano somministrati il PANAS post, il Godspeed e il questionario MDMT adattato, seguiti da un breve debriefing finale in cui raccoglievo anche impressioni più libere sulla sessione.

Nel protocollo la finestra di interazione era stata pensata in modo controllato, ma nella pratica la durata complessiva dipendeva anche dai passaggi di setup. Per questo, nelle prove reali, la sessione si è assestata in media attorno ai 25 minuti nella condizione neutra e ai 30 minuti nella condizione personale. La Figura~\ref{fig:empirical-procedure} riassume il flusso seguito durante il test.

\begin{figure}[H]
\centering
\includegraphics[width=0.60\textwidth]{empirical_procedure.png}
\caption[Procedura sperimentale]{Sequenza della procedura sperimentale usata nello studio empirico.}
\label{fig:empirical-procedure}
\end{figure}

\section{Risultati}
\label{sec:empirical-results}

\paragraph{Misure di base.}

I risultati vanno letti con prudenza. Per i confronti tra gruppi ho usato Mann--Whitney U, mentre per le variazioni pre/post dentro la stessa condizione ho usato Wilcoxon per ranghi con segno. In un campione piccolo come questo, p-value ed effetto servono soprattutto a leggere la direzione dei risultati e a capire dove il prototipo mostra segnali coerenti.

Prima degli outcome principali, la Tabella~\ref{tab:empirical-baseline} mostra che autostima e affettività positiva pre-test risultano abbastanza vicine, mentre l'affettività negativa iniziale parte più alta nella condizione personale. Il confronto resta leggibile, ma questa differenza iniziale va tenuta presente quando si commentano i cambiamenti emotivi.

Nelle tabelle che seguono, PA indica l'affettività positiva (\textit{Positive Affect}), NA l'affettività negativa (\textit{Negative Affect}) e il simbolo $\Delta$ la variazione tra post-test e pre-test.

\begin{table}[H]
\centering
\small
\setlength{\tabcolsep}{5pt}
\renewcommand{\arraystretch}{1.05}
\begin{tabularx}{0.94\textwidth}{>{\raggedright\arraybackslash}X>{\raggedleft\arraybackslash}X>{\raggedleft\arraybackslash}X>{\raggedleft\arraybackslash}p{0.09\textwidth}>{\raggedleft\arraybackslash}p{0.09\textwidth}}
\hline
\textbf{Misura} & \textbf{Personale mediana [IQR]} & \textbf{Neutra mediana [IQR]} & \textbf{p} & \textbf{r\textsubscript{rb}} \\
\hline
RSES totale & 30.00 [26.25--31.75] & 32.00 [29.25--35.25] & 0.150 & -0.31 \\
PANAS pre -- PA & 20.00 [16.50--21.75] & 18.50 [17.00--20.50] & 0.517 & 0.14 \\
PANAS pre -- NA & 9.50 [8.00--13.75] & 7.00 [6.00--8.00] & 0.008 & 0.57 \\
\hline
\end{tabularx}
\caption{Confronto descrittivo delle misure di base prima dell'interazione.}
\label{tab:empirical-baseline}
\end{table}

\begin{table}[H]
\centering
\small
\setlength{\tabcolsep}{5pt}
\renewcommand{\arraystretch}{1.05}
\begin{tabularx}{0.94\textwidth}{>{\raggedright\arraybackslash}X>{\raggedleft\arraybackslash}X>{\raggedleft\arraybackslash}X>{\raggedleft\arraybackslash}p{0.09\textwidth}>{\raggedleft\arraybackslash}p{0.09\textwidth}}
\hline
\textbf{Outcome} & \textbf{Personale mediana [IQR]} & \textbf{Neutra mediana [IQR]} & \textbf{p} & \textbf{r\textsubscript{rb}} \\
\hline
MDMT totale & 4.50 [3.40--5.67] & 4.55 [3.84--5.39] & 0.983 & 0.01 \\
$\Delta PA$ & 0.00 [-1.00--1.00] & 1.00 [0.00--2.00] & 0.105 & -0.35 \\
$\Delta NA$ & -2.00 [-3.75--0.00] & 0.00 [-1.25--0.00] & 0.165 & -0.29 \\
\hline
\end{tabularx}
\caption[Outcome principali nelle due condizioni]{Confronto tra le due condizioni sugli outcome principali dello studio empirico.}
\label{tab:empirical-core-outcomes}
\end{table}

\begingroup
\small
\setlength{\tabcolsep}{5pt}
\renewcommand{\arraystretch}{1.05}
\begin{longtable}{>{\raggedright\arraybackslash}p{0.26\textwidth}>{\raggedleft\arraybackslash}p{0.20\textwidth}>{\raggedleft\arraybackslash}p{0.20\textwidth}>{\raggedleft\arraybackslash}p{0.09\textwidth}>{\raggedleft\arraybackslash}p{0.09\textwidth}}
\caption[Dimensioni percettive dell'avatar nelle due condizioni]{Confronto tra le due condizioni sulle principali dimensioni percettive dell'avatar.}
\label{tab:empirical-godspeed}\\
\hline
\textbf{Dimensione} & \textbf{Personale mediana [IQR]} & \textbf{Neutra mediana [IQR]} & \textbf{p} & \textbf{r\textsubscript{rb}} \\
\hline
\endfirsthead
\hline
\textbf{Dimensione} & \textbf{Personale mediana [IQR]} & \textbf{Neutra mediana [IQR]} & \textbf{p} & \textbf{r\textsubscript{rb}} \\
\hline
\endhead
\hline
\endfoot
Antropomorfismo & 3.00 [2.00--3.50] & 2.60 [2.40--3.05] & 0.573 & 0.12 \\
Animazione & 3.12 [2.50--3.50] & 2.88 [2.19--3.31] & 0.544 & 0.13 \\
Gradevolezza & 3.80 [3.45--3.95] & 3.80 [2.95--4.40] & 0.967 & -0.01 \\
Intelligenza percepita & 3.30 [2.70--4.00] & 3.10 [2.15--3.70] & 0.358 & 0.20 \\
Sicurezza percepita & 4.33 [3.42--5.00] & 4.33 [3.67--4.67] & 0.751 & 0.07 \\
\hline
\end{longtable}
\endgroup

\subsection{Fiducia e percezione dell'avatar}
\label{subsec:trust-avatar}

Il dato più chiaro è che la personalizzazione, da sola, non ha prodotto un aumento evidente della fiducia dichiarata: il punteggio MDMT resta quasi sovrapponibile nelle due condizioni. Almeno nello stato attuale del prototipo, volto, voce e memoria iniziale non bastano quindi a far percepire l'agente come più affidabile.

Sulla percezione dell'avatar emerge qualcosa in più. La condizione personale tende a stare leggermente sopra la neutra soprattutto nel modo in cui l'agente viene percepito come più umano, un po' più animato e leggermente più credibile sul piano cognitivo, mentre la gradevolezza resta sostanzialmente allineata. Non è una differenza forte, ma va nella direzione attesa: l'avatar personale viene letto un po' di più come un interlocutore riconoscibile e meno come una figura generica.

\begin{figure}[H]
\centering
\includegraphics[width=0.93\textwidth]{empirical_godspeed_bars.png}
\caption[Dimensioni percettive dell'avatar]{Confronto delle principali dimensioni percettive dell'avatar nelle due condizioni; nel grafico la barra chiara indica la condizione personale, quella scura la neutra.}
\label{fig:empirical-godspeed}
\end{figure}

Qui la lettura deve restare indiretta: nel materiale raccolto non c'è una misura specifica dell'identificazione con l'avatar, quindi non posso dire che la condizione personale la aumenti in senso forte. Posso però osservare che alcuni indicatori percettivi si muovono nella direzione di una maggiore riconoscibilità dell'agente.

\subsection{Variazione dello stato affettivo}
\label{subsec:affective-change}

Per la parte affettiva ho guardato soprattutto a $\Delta PA$ e $\Delta NA$, cioè alla differenza tra post-test e pre-test nelle due componenti del PANAS. I confronti tra gruppi non arrivano a una differenza statisticamente netta, ma mostrano due andamenti diversi: nella condizione neutra l'affettività positiva tende a crescere un po' di più, mentre nella condizione personale si vede una riduzione maggiore dell'affettività negativa.

Anche leggendo i test entro ciascun gruppo, il dato più interessante resta nella condizione personale, dove la componente negativa si abbassa in modo più visibile tra inizio e fine sessione ($p = 0.022$; $r_{rb} = -0.73$). Nella condizione neutra, invece, si nota soprattutto una tendenza all'aumento dell'affettività positiva ($p = 0.057$; $r_{rb} = 0.57$), che resta però appena sotto una soglia di solito letta come più convincente.

Il dato, però, non va forzato oltre quello che mostra: il campione resta piccolo, i gruppi non partono dallo stesso livello di affettività negativa e non c'è nessuna base per parlare di efficacia psicologica del sistema. L'impressione più netta è che l'avatar personale incida soprattutto sul tono della sessione, più che sulla fiducia dichiarata in modo esplicito.

\subsection{Analisi qualitativa dei log}
\label{subsec:qualitative-logs}

I questionari da soli non bastavano a spiegare alcuni comportamenti osservati durante il test. Per questo ho affiancato una rilettura qualitativa dei log conversazionali e delle note prese durante le sessioni. Ho usato un codebook volutamente breve, costruito attorno ai temi che tornavano più spesso:

\begingroup
\small
\setlength{\tabcolsep}{6pt}
\renewcommand{\arraystretch}{1.05}
\begin{longtable}{>{\raggedright\arraybackslash}p{0.30\textwidth}>{\raggedright\arraybackslash}p{0.62\textwidth}}
\caption{Codebook usato per la lettura qualitativa delle note di test.}
\label{tab:qualitative-codebook}\\
\hline
\textbf{Categoria} & \textbf{Che cosa intercetta nelle note} \\
\hline
\endfirsthead
\hline
\textbf{Categoria} & \textbf{Che cosa intercetta nelle note} \\
\hline
\endhead
\hline
\endfoot
Richiamo corretto della memoria & recupero di un'informazione personale o contestuale in modo coerente con il profilo \\
Dettaglio inventato / allucinazione & aggiunta di informazioni non presenti nella memoria o nel turno corrente \\
Ripetitività & risposte che tornano sugli stessi fatti o sugli stessi aggettivi senza avanzare davvero \\
Resa migliore con domande specifiche & casi in cui richieste più precise producono risposte più centrate \\
Coinvolgimento / divertimento & segnali di curiosità, sorpresa o interesse emersi durante la prova \\
Osservazioni su voce e interfaccia & commenti su voce sintetica, orientamento nella UI e resa generale del sistema \\
\hline
\end{longtable}
\endgroup

\begin{figure}[H]
\centering
\includegraphics[width=0.93\textwidth]{empirical_log_themes.png}
\caption[Temi ricorrenti nei log e nelle osservazioni]{Temi ricorrenti emersi nelle note di osservazione e nei log conversazionali; la barra chiara indica la condizione personale, quella scura la neutra.}
\label{fig:empirical-log-themes}
\end{figure}

I conteggi della Figura~\ref{fig:empirical-log-themes} non vanno letti come misura turno per turno, ma come presenza del tema almeno una volta nelle note legate a un partecipante. La differenza più netta è che la condizione personale raccoglie più casi di richiamo corretto della memoria, ma anche più casi di dettaglio inventato. È un esito plausibile: appena il sistema riceve più informazioni personali, ha anche più occasioni per funzionare bene e più punti in cui può sbagliare in modo evidente.

Letti in modo aggregato, questi pattern aiutano a spiegare meglio anche i questionari. Nella condizione personale il sistema produce più occasioni in cui il richiamo della memoria viene percepito come centrato, e questo rende l'interazione più riconoscibile quando il recupero funziona. Allo stesso tempo, però, è anche la condizione in cui gli errori diventano più visibili: un dettaglio inventato su hobby, città o caratteristiche personali pesa più di una risposta generica data da un avatar neutro.

La condizione neutra mostra un profilo diverso. Le sessioni risultano in media meno specifiche, ma anche meno esposte a rotture nette di credibilità legate al profilo personale. Nelle note di osservazione torna più spesso un'interazione semplice da sostenere, a tratti persino più leggera, mentre nella condizione personale il coinvolgimento dipendeva molto di più da quanto il sistema riuscisse davvero a usare bene la memoria iniziale.

Tre mini-casi aiutano a rendere più concreto questo passaggio. \textbf{Caso 1.} In una sessione della condizione personale, una domanda sociale poco ancorata al profilo ha portato prima a una risposta troppo sicura su una terza persona e subito dopo a un ripiegamento difensivo del tipo ``non ho abbastanza elementi'' quando l'utente ha contestato il contenuto. Qui il problema non era solo l'errore, ma il fatto che l'avatar passasse in pochi turni da un'invenzione piuttosto libera a una chiusura brusca. È un esempio chiaro di come la generalizzazione possa pesare direttamente sulla credibilità dell'interazione.

\textbf{Caso 2.} In un'altra sessione, sempre personale, una domanda iniziale su un obiettivo di vita riceveva una risposta prudente perché il sistema non aveva ancora un riferimento abbastanza preciso. Dopo che l'utente aveva aggiunto quel dettaglio in modo esplicito, una domanda successiva sullo stesso tema produceva invece una risposta molto più centrata e coerente con il profilo. Quando l'informazione viene formulata in modo specifico, il retrieval riesce più facilmente a recuperarla e a usarla in modo utile.

\textbf{Caso 3.} In una sessione più leggera, con domande semplici su gusti, attività quotidiane e piccoli ricordi, lo scambio è rimasto fluido per più turni e l'avatar ha mantenuto un tono abbastanza scorrevole. Il limite, però, riemergeva appena l'utente provava a far ricordare qualcosa di nuovo su di sé: il sistema recuperava una parte corretta dell'informazione, ma la mescolava con dettagli appartenenti al profilo dell'avatar o con elementi troppo generici. Il caso è utile perché mostra insieme i due lati più frequenti del test: interazione facile da sostenere e coinvolgimento discreto, ma robustezza ancora incostante quando il contesto cambia.

Il fatto che l'MDMT rimanga quasi sovrapposto nelle due condizioni suggerisce che la fiducia dichiarata dipenda ancora soprattutto dalla qualità concreta della risposta, non dalla sola personalizzazione dell'avatar. In parallelo, la riduzione di affettività negativa osservata nella condizione personale è compatibile con le sessioni in cui il richiamo corretto della memoria ha dato all'interazione un tono più vicino; l'andamento di $\Delta PA$ nella condizione neutra, invece, si accorda meglio con sessioni percepite come meno dense sul piano personale, ma più facili da sostenere fino alla fine.

Anche sul piano tecnico il quadro resta abbastanza coerente. Quando la domanda è specifica, il recupero della memoria riesce più spesso a fornire un contesto utile alla risposta; quando invece la richiesta resta vaga, il sistema tende più facilmente a generalizzare oppure a introdurre dettagli non sufficientemente ancorati ai dati disponibili. Per questo i log non vengono usati qui come aneddoti, ma come supporto interpretativo ai risultati emersi dai questionari.

\section{Discussione}
\label{sec:empirical-discussion}

Lo studio empirico non mostra un vantaggio netto della condizione personale su tutti i fronti, ma chiarisce abbastanza bene dove la personalizzazione comincia a cambiare l'interazione. Sul piano della fiducia dichiarata il risultato resta quasi piatto: l'MDMT è molto simile tra i due gruppi, segno che la fiducia dipende ancora soprattutto dalla qualità concreta della risposta, non dalla sola somiglianza dell'avatar con la persona.

Sul piano percettivo il quadro è più sfumato. L'avatar personale tende a risultare un po' più riconoscibile, leggermente più vivo nella resa e un po' meno generico, ma il guadagno resta contenuto. La personalizzazione aiuta a rendere l'incontro più distinguibile, ma non basta a compensare i punti in cui retrieval e generazione restano fragili \cite{Yang2025ECA_XR,Kan2023ECA_Training}.

La parte affettiva è forse quella più interessante. La riduzione dell'affettività negativa nella condizione personale non autorizza letture forti, ma lascia intravedere che riconoscere qualcosa di proprio nell'avatar possa incidere sul tono della sessione. Il gruppo neutro, invece, sembra muoversi in modo un po' diverso: meno riconoscibilità, ma in alcuni casi un'interazione vissuta come più leggera o più ludica.

La lettura qualitativa aiuta a capire meglio il perché. Nella condizione personale il sistema riesce più spesso a recuperare dettagli giusti e a produrre momenti percepiti come ``vicini'', ma è anche la condizione in cui gli errori pesano di più. Un dettaglio inventato su hobby, città o carattere rompe molto più facilmente la credibilità dell'interazione rispetto a quanto accade con un avatar neutro.

Restano comunque limiti metodologici che è meglio tenere espliciti. Il campione è ridotto, il confronto è tra gruppi, i gruppi non partono dallo stesso livello di affettività negativa e manca una misura diretta dell'identificazione con l'avatar. Anche l'analisi qualitativa viene da note e log riletti a posteriori, quindi funziona più come triangolazione dei risultati che come misura autonoma.

Questo test aggiunge qualcosa che la sola valutazione tecnica non poteva dare: il prototipo non viene percepito solo come una pipeline che funziona oppure no, ma come un interlocutore che in alcuni momenti riesce davvero a dare un minimo di continuità personale. Chiarisce anche dove il sistema si rompe: non tanto sulla sola latenza, quanto nel momento in cui deve capire quando usare la memoria, quanto fidarsi del contesto recuperato e quando invece restare più prudente.
